\documentclass[10pt,a4paper]{amsart}
\usepackage[margin=19mm]{geometry}
\usepackage{amsmath,amssymb,mathtools}
\usepackage[scr=rsfs,cal=euler]{mathalfa}
\usepackage{xfrac}
\usepackage[unicode,hidelinks,bookmarks=false]{hyperref}
\newcommand{\ie}{i.e.\ }
\newcommand{\cf}{cf.\ }
\newcommand{\resp}{respectively\ }
\newcommand{\Subsection}{\subsection}
\providecommand{\nobreakdash}{\nobreak\hbox{-}}
\setlength{\emergencystretch}{2em}
\multlinegap0pt
\usepackage{bm}
\usepackage{bbm}
\usepackage{dsfont}
\usepackage{xspace}
\usepackage{cancel}
\usepackage{mathtools}
\usepackage{cleveref}
\usepackage[shortlabels]{enumitem}
\usepackage{amsthm}
\makeatletter
\@ifundefined{dummy}{\newtheorem{dummy}{Dummy}}{}
\@ifundefined{lemma}{\newtheorem{lemma}[dummy]{Lemma}}{}
\@ifundefined{theorem}{\newtheorem{theorem}[dummy]{Theorem}}{}
\makeatother
\DeclareMathOperator{\high}{\mathrm{high}}
\DeclareMathOperator{\low}{\mathrm{low}}
\newcommand{\E}{\mathbb{E}}
\newcommand{\calC}{\mathcal{C}}
\newcommand{\calF}{\mathcal{F}}
\newcommand{\calH}{\mathcal{H}}
\newcommand{\ignore}[1]{{}}
\newcommand{\poly}{\mathsf{poly}}
\newcommand{\shachar}[1]{\textcolor{blue}{(Shachar: {#1})}}
\newcommand{\supp}{\mathrm{supp}}
\newcommand{\yimeng}[1]{\textcolor{orange}{(Yimeng: {#1})}}
\renewcommand{\Pr}{\mathop{\bf Pr\/}}
\renewcommand{\epsilon}{\varepsilon}
\title[Moonflowers and efficient code sparsification]{Moonflowers and efficient code sparsification}
\author{Shachar Lovett, Raghu Meka, Yimeng Wang}
\date{Source: arXiv:2605.08676v2 (CC BY 4.0)}
\begin{document}
\maketitle

\noindent\emph{Source and licence.} Moonflowers and efficient code sparsification. Version arXiv:2605.08676v2 (CC BY 4.0), distributed under the Creative Commons Attribution 4.0 International licence, as recorded in the arXiv licence metadata for this exact version and shown on the abstract page https://arxiv.org/abs/2605.08676v2. Full rights notice: licenses/arxiv-2605.08676-CC-BY-4.0.txt.\par\smallskip

\noindent\textit{Editorial scope.} Counted unit: the Theorem whose source label is \texttt{thm:sparsif\_single\_log\_n\_v2} in the article (source0.tex lines 1359--1370, source label \texttt{thm:sparsif\_single\_log\_n\_v2}), delivered together with the article's own complete original proof of it (source0.tex lines 1431--1779, one \texttt{proof} environment of 349 lines). No mathematics is written, repaired or re-derived: statement and proof are the article's own text line for line. The proof is detached from the statement in the source: the article states the result at lines 1359--1370 and prints its proof at lines 1431--1779, and the proof environment's own header names the statement, so the association is the article's own. Required context retained uncounted: main\_wsf\_theorem (lines 380--384); thm:universe\_reduction (lines 417--419); lem:WSF\_free\_preserved\_under\_projection (lines 584--585); lem:moonflower\_free\_family\_have\_small\_support (lines 604--609); lem:sparsif\_chernoff\_v2 (lines 720--725). Every definition, display and label of those lines is retained unchanged. Omitted from this package and not redistributed: 1--379, 385--416, 420--583, 586--603, 610--719, 726--1358, 1371--1430, 1780--1894. Nothing inside a retained excerpt is omitted or altered: every retained body is delivered exactly as the article's lines are written, so no omission receipt of any kind accompanies this package. Citations: the retained excerpts carry no citation command at all, so no bibliography entry of the article is transcribed and no reference is redistributed. Reference adaptation: none was needed and none was made; every reference inside the delivered unit resolves inside it, so no unresolved reference or citation is produced. Printed slips kept literal: no printed word, symbol, hypothesis or number of the counted statement or of its proof was altered. Layout and class: the article's own class is \texttt{article} with packages including graphicx, amsmath, amssymb, amsthm, fullpage, fontenc, geometry, latexsym, caption2, graphics, epsfig, amsopn, url, xcolor; the wrapper is the lane's pinned \texttt{amsart} editorial wrapper at 10pt on A4, which restates the theorem-like environments the retained text needs and adds the article's own macro definitions that the retained text needs (\texttt{\textbackslash E}, \texttt{\textbackslash Pr}, \texttt{\textbackslash calC}, \texttt{\textbackslash calF}, \texttt{\textbackslash calH}, \texttt{\textbackslash epsilon}, \texttt{\textbackslash high}, \texttt{\textbackslash ignore}, \texttt{\textbackslash low}, \texttt{\textbackslash poly}, \texttt{\textbackslash shachar}, \texttt{\textbackslash supp}, \texttt{\textbackslash yimeng}), with extra wrapper packages (bm, bbm, dsfont, xspace, cancel, mathtools, cleveref, enumitem). The delivered numbering is the unit's own. Assets: no retained span contains a figure environment, a graphics-inclusion command, a PDF-page inclusion, a table or a TikZ picture; no image file is copied into this package, the package declares no asset dependency and no image file is redistributed with it. Licence: version arXiv:2605.08676v2 of this article is distributed under the Creative Commons Attribution 4.0 International licence, as recorded in the arXiv licence metadata for this exact version and shown on the abstract page https://arxiv.org/abs/2605.08676v2. A full rights notice is delivered at licenses/arxiv-2605.08676-CC-BY-4.0.txt. Characters: every retained excerpt is pure ASCII, so the delivered engine, XeLaTeX with its default Latin Modern text font, reports no missing character.
\begin{theorem}[Extremal bounds on moonflowers]\label{main_wsf_theorem}  Let $\calF$ be a family of $w$-sets which is $k$-moonflower-free. Then 
\[
|\calF| \leq \binom{k+w-1}{w}.
\]
\end{theorem}

\begin{theorem}[Universe reduction]\label{thm:universe_reduction}
    Let $\calF \subseteq 2^{[n]}$ be a family of $w$-sets which is $k$-moonflower-free. Then for any $\lambda \in (0, 1]$, there exists $I \subseteq [n]$ of size $|I| \leq O(\lambda^{-1}k \log^3 (kw / \lambda))$ satisfying $|\calF_{\bar{I}}| \leq \exp (\lambda w)$.
\end{theorem}

\begin{lemma}[Moonflower-freeness is preserved under projection]\label{lem:WSF_free_preserved_under_projection} If $\calF \subseteq 2^{[n]}$ is $k$-moonflower-free and $J \subseteq [n]$, then $\calF_J$ is $k$-moonflower-free.
\end{lemma}

\begin{lemma}[Support bound from moonflower-freeness]\label{lem:moonflower_free_family_have_small_support}Let $\calF \subseteq 2^{[n]}$ be a family of $w$-sets which is $k$-moonflower free. Then
\[
|\supp (\calF)| \leq (k-1)w.
\]
In particular, after deleting unused coordinates, we may assume $n \leq (k-1)w$. 
\end{lemma}

\begin{lemma}[Chernoff bound, convenient form]\label{lem:sparsif_chernoff_v2}
Let $t\ge 1$ and let $X\sim \mathrm{Bin}(t,1/2)$. Then for every $\Delta>0$,
\[
\Pr\left[\,|2X-t|>\Delta\,\right]\ \le\ 2\exp\!\left(-\frac{\Delta^2}{3t}\right).
\]
\end{lemma}

\begin{theorem}[Improved sparsification with a single $\log n$ factor]
\label{thm:sparsif_single_log_n_v2}
Let $\calC\subseteq\{0,1\}^n$ satisfy $\mathrm{NRD}(\calC)\le k-1$.
Then for every $\varepsilon\in(0,1/4)$ there exists a weighted coordinate set $(T,\alpha)$
that $\varepsilon$-sparsifies $\calC$ and satisfies
\[
|T|
\ \le\
\frac{k\log n}{\varepsilon^2}\cdot \poly\!\left(\log(k/\varepsilon),\log\log n\right).
\]
%Moreover, $(T,\alpha)$ can be obtained by the randomized recursive construction below with success probability at least $2/3$.
\end{theorem}

\begin{proof}(of \Cref{thm:sparsif_single_log_n_v2})

\paragraph{Parameter choices.}
Fix $R:=\lceil \log_2 n\rceil$. In what follows, let $C$ be a large constant. Define the transition threshold \ignore{\shachar{we have 3 unspecified constants $C$ below, do we assume they are all the same? if not then lets call them $C_1,C_2,C_3$; this then needs to be propagated into the proof} \yimeng{I think we can assume they are the same. As long as it's large enough, the proof should go through.}\shachar{at least it should be clear if $C$ should be large enough or small enough. EG if we incrase $C$ does it keep everything consistent?}}
\[
w_{\mathrm{high}}\ :=\ \left\lceil
C\cdot \frac{k\log(k/\varepsilon)+\log R}{\varepsilon^2}
\right\rceil.
\]
Define the medium-weight per-round error
\[
\eta_0\ :=\ \frac{\varepsilon}{100\log(2w_{\mathrm{high}})}.
\]
Define the low-weight cutoff
\[
w_{\mathrm{low}}\ :=\ \left\lceil \frac{C}{\eta_0^2}\cdot\left(\log(k/\varepsilon)+\log R\right)\right\rceil.
\]
Finally, for dyadic $w>w_{\mathrm{high}}$, define the large-weight per-round error
\[
\eta(w)\ :=\ \min\left\{\frac14,\ \sqrt{\frac{C\,(k\log( w/k)+\log R)}{w}}\right\}.
\]

\paragraph{Recursive construction.}
Let $U_0=[n]$. For rounds $r=0,1,\dots,R-1$, given $U_r$ we choose a puncturing set $I_r\subseteq U_r$,
set $V_r:=U_r\setminus I_r$, and then form $U_{r+1}$ by including each element of $V_r$
independently with probability $1/2$.

For dyadic\footnote{Here by dyadic, we mean $w = 2^j w_{\low}$ for $j = 0,1,...,\log (w_{\high} / w_{\low})$.} $w\in[w_{\low},w_{\high}]$, define
\[
\calF^{(r)}_{(w,2w]}
:=
\left\{\,\supp(x)\cap U_r:\ x\in\calC,\ \ w<|\supp(x)\cap U_r|\le 2w\,\right\}.
\]
Define the low-weight family
\[
\calF^{(r)}_{\le w_{\low}}
:=
\left\{\,\supp(x)\cap U_r:\ x\in\calC,\ \ |\supp(x)\cap U_r|\le w_{\low}\,\right\}.
\]
(All these families are $k$-moonflower-free by \Cref{lem:WSF_free_preserved_under_projection}.)

\smallskip
\begin{enumerate}[(a)]
    \item \emph{Low weights}: capture deterministically. Let
    \[
    I_{r,\mathrm{low}}:=\supp\!\left(\calF^{(r)}_{\le w_{\low}}\right)\subseteq U_r.
    \]
    By \Cref{lem:moonflower_free_family_have_small_support}, $|I_{r,\mathrm{low}}|\le k\,w_{\low}$.
    \item \emph{Medium weights}: puncture to shrink traces. Choose a sufficiently small absolute constant $\theta_0 > 0$, say $\theta_0 = 10^{-3}$ after adjusting constants in
    the Chernoff bound. 
    For each dyadic $w\in[w_{\low},w_{\high}]$, apply \Cref{thm:universe_reduction}
    to $\calF^{(r)}_{(w,2w]}$ with parameter $\lambda = \theta_0 \eta_0^2$ to obtain $I_{r,w}\subseteq U_r$
    such that
    \begin{equation}\label{eq:medium_trace_bound_v2}
    \big|\left(\calF^{(r)}_{(w,2w]}\right)_{\overline{I_{r,w}}}\big|
    \ \le\
    \exp (\theta_0 \eta_0^2 w),
    \end{equation}
    and
    \[
    |I_{r, w}| \leq \frac{k}{\eta_0^2} \cdot \poly \left(\log \frac{kw}{\eta_0} \right).
    \]
    Here we use the fact that $\eta_0 < 1$ and $\theta_0$ is an absolute constant, so its contribution is absorbed into the polylogarithmic factor.
    \item \emph{Define $I_r$ and recurse.} Let
    \[
    I_r\ :=\ I_{r,\mathrm{low}} \ \cup\ \bigcup_{\substack{\text{dyadic }w\\ w_{\low}\le w\le w_{\high}}} I_{r,w},
    \qquad
    V_r:=U_r\setminus I_r,
    \]
    and sample $U_{r+1}\subseteq V_r$ by keeping each coordinate with probability $1/2$.
    \item \emph{Output.} Define
    \[
    T\ :=\ \left(\bigcup_{r=0}^{R-1} I_r\right)\ \cup\ U_R,
    \qquad
    \alpha(i):=2^r\ \text{ if } i\in I_r,\quad \alpha(i):=2^R\ \text{ if } i\in U_R.
    \]
    (The sets $I_0,\dots,I_{R-1},U_R$ are disjoint since $U_{r+1}\subseteq U_r\setminus I_r$.)
\end{enumerate}



\paragraph{Successful Sampling.}
Fix a round $r$ and condition on the entire history up to round $r$ so that $U_r,I_r,V_r$ are fixed.

\smallskip
\emph{Medium regime.}
Fix dyadic $w\in[w_{\low},w_{\high}]$. Consider the trace family on $V_r$:
\[
\left(\calF^{(r)}_{(w,2w]}\right)_{V_r}.
\]
Since $V_r\subseteq U_r\setminus I_{r,w}$, restriction cannot increase size, so
\[
\big|\left(\calF^{(r)}_{(w,2w]}\right)_{V_r}\big|
\le
\big|\left(\calF^{(r)}_{(w,2w]}\right)_{\overline{I_{r,w}}}\big| \leq \exp (\theta_0 \eta_0^2 w).
\]
Moreover, for any $A$ in this family, we have $|A|\le 2w$.
Applying \Cref{lem:sparsif_chernoff_v2} with $\Delta:=\eta_0 w$ gives 
\[
\Pr_{U_{r+1}}\left[\big|\,2|A\cap U_{r+1}|-|A|\,\big|>\eta_0 w\right]
\le
2\exp\!\left(-\Omega(\eta_0^2 w)\right).
\]
Using~\eqref{eq:medium_trace_bound_v2} and the definition of $w_{\low}$,
the union bound over all $A\in (\calF^{(r)}_{(w,2w]})_{V_r}$ fails with probability at most
$\exp(-\Omega(\eta_0^2 w))\le 1/(100R\log(2w_{\high}))$ for appropriate choice of constants $C, \theta_0$.

\smallskip
\emph{Large regime.}
Fix dyadic $w>w_{\high}$ and consider the trace family
\[
\calH^{(r)}_{(w,2w]}\ :=\ \left(\calF^{(r)}_{(w,2w]}\right)_{V_r}.
\]
Since $\calF^{(r)}_{(w,2w]}$ is $k$-moonflower-free and $w\ge k$ in this regime,
\Cref{main_wsf_theorem} gives 
\[
|\calH^{(r)}_{(w,2w]}|
\ \le\
|\calF^{(r)}_{(w,2w]}|
\ \le\
\left(\frac{C_0 w}{k}\right)^k,
\]
for some constant $C_0 > 0$.
For any $A\in\calH^{(r)}_{(w,2w]}$ we have $|A|\le 2w$, so by \Cref{lem:sparsif_chernoff_v2} with
$\Delta:=\eta(w)\,w$, 
\[
\Pr_{U_{r+1}}\left[\big|\,2|A\cap U_{r+1}|-|A|\,\big|>\eta(w)\,w\right]
\le
2\exp\!\left(-\Omega(\eta(w)^2 w)\right)
=
2\exp\!\left(-\Omega(k\log(e w/k)+\log R)\right).
\]
For large enough $C$ in the definition of $\eta(w)$, the union bound over all
$A\in\calH^{(r)}_{(w,2w]}$ fails with probability at most $1/(100R\cdot 2^{j+2})$
when $w=2^j$. Summing over dyadic $w>w_{\high}$ gives failure probability at most $1/(100R)$.

\smallskip
Combining medium and large regimes and summing over the $O(\log w_{\high})$ medium scales,
we obtain that conditioned on the past, round $r$ fails with probability at most $1/(50R)$.
A union bound over $r=0,\dots,R-1$ yields that all rounds succeed with probability at least $49/50$.
Also $\E[|U_R|]\le n2^{-R}\le 1$, hence $\Pr[|U_R|\le 50]\ge 49/50$ by Markov.
Intersecting gives overall success probability at least $2/3$.

\paragraph{Accuracy for a fixed codeword.}
Fix a successful outcome as above. Fix a codeword $x\in\calC$ with $S:=\supp(x)$.

Let $\widehat{N}_r = \sum_{i=0}^{r-1} 2^i |S \cap I_i| + 2^r |S \cap U_r|$ denote the weight estimate for the weight of $x$ at round $r$. Note that $\widehat{N}_0 = |S|$ is the Hamming weight of $x$. Our goal is to show that $\widehat{N}_R = (1\pm \varepsilon) \widehat{N}_0$. 

Consider the next weight estimate  $$\widehat{N}_{r+1} = \sum_{i=0}^{r-1} 2^i  |S \cap I_i| + 2^r (|S \cap I_r| + 2 |S \cap U_{r+1}|).$$

Let $w_r := |S \cap U_r|$ be the true residual weight at round $r$, and let 

$$
\epsilon_r := \begin{cases} 0 & w_r \leq w_{\low},\\
\eta_0 & w_{\low} < w_r \leq w_{\high}\\
\eta(w_r) & w_r > w_{\high}
\end{cases}.
$$

Then, as we are in a successful outcome, we have,
$$\left|2 |S \cap U_{r+1}| - |S \cap (U_r\setminus I_r)| \right| \leq \epsilon_r w_r.$$

Thus, 
$$\left| 2 |S \cap U_{r+1}| + |S \cap I_r|  - |S \cap U_r| \,\right| \leq \epsilon_r w_r.$$

Combining the above equations, we get 
$$|\widehat{N}_{r+1} - \widehat{N}_r| = 2^r \cdot \left| 2 |S \cap U_{r+1}| + |S \cap I_r|  - |S \cap U_r| \,\right| \leq \epsilon_r 2^r w_r \leq \epsilon_r \widehat{N}_r.$$

Thus, we have 
$$(1-\epsilon_r) \widehat{N}_r \leq \widehat{N}_{r+1} \leq (1+\epsilon_r) \widehat{N_r}.$$

Applying the above inequality for $r = 0,\ldots,R-1$, we get 
$$\widehat{N}_0  \cdot \prod_{r=0}^{R-1} (1-\epsilon_r) \leq \widehat{N}_R \leq \widehat{N}_0  \cdot \prod_{r=0}^{R-1} (1+\epsilon_r).$$

Further, as $\epsilon_r < 1/2$, the above can be simplified to 
\begin{equation}\label{eq:sparse1}
    \widehat{N}_0  \cdot \exp\left(-2 \sum_{r=0}^{R-1}\epsilon_r\right) \leq \widehat{N}_R \leq \widehat{N}_0  \cdot \exp\left( \sum_{r=0}^{R-1}\epsilon_r\right) .
\end{equation}

It remains to bound $\sum_r \epsilon_r$. The number of rounds with $w_{\low}<w_r\le w_{\high}$ is $O(\log(2w_{\high}))$ since
$w_{r+1}\le (1+\eta_0)w_r/2\le (3/5)w_r$.
Thus $\sum_{w_{\low}<w_r\le w_{\high}}\epsilon_r\le O(\eta_0\log(2w_{\high}))\le \varepsilon/50$.

For rounds with $w_r>w_{\high}$, we have $\epsilon_r\le \eta(w_r)$ and 
$w_{r+1}\le (1+\epsilon_r)w_r/2\le (5/8)w_r$. A geometric-series estimate gives
\[
\sum_{w_r>w_{\high}}\epsilon_r
\ \le\
O\!\left(\sqrt{\frac{k\log(e w_{\high}/k)+\log R}{w_{\high}}}\right)
\ \le\ \varepsilon/50
\]
by the definition of $w_{\high}$ (with $C$ large enough).

Combining the above we get that $\sum_r \epsilon_r \leq \varepsilon/25$. Thus, in particular we must have 
$\widehat{N}_R = (1 \pm \varepsilon) N_0$, as we wanted.

\ignore{
Let $W_r:=|S\cap U_r|$ be the true residual weight at round $r$.  
Define the estimator recursively:
\[
\widehat W_R:=|S\cap U_R|,
\qquad
\widehat W_r:=|S\cap I_r|+2\widehat W_{r+1}\quad (r=R-1,\dots,0).
\]
Then $\widehat W_0=\widehat{\mathrm{wt}}_{T,\alpha}(x)$.

Let $R_r:=S\cap V_r$, so $W_r=|S\cap I_r|+|R_r|$ and $W_{r+1}=|R_r\cap U_{r+1}|$.
If $W_r\le w_{\min}$ then $S\cap U_r\subseteq I_{r,\mathrm{tiny}}\subseteq I_r$, so $R_r=\emptyset$
and $W_{r+1}=0$ and $\widehat W_r=W_r$ exactly.

Otherwise, let dyadic $w$ satisfy $W_r\in(w,2w]$.
If $w\le w_{\high}$ then by the medium-regime guarantee we have
\[
\big|\,2W_{r+1}-|R_r|\,\big|\ \le\ \eta_0 w\ \le\ \eta_0 W_r.
\]
If $w>w_{\high}$ then by the large-regime guarantee we have
\[
\big|\,2W_{r+1}-|R_r|\,\big|\ \le\ \eta(w)\,w\ \le\ \eta(w)\,W_r.
\]
In either case we obtain
\begin{equation}\label{eq:eta_r_def_v2}
\big|\,2W_{r+1}-|R_r|\,\big|\ \le\ \eta_r\,W_r,
\qquad\text{where}\qquad
\eta_r:=\begin{cases}
0 & W_r\le w_{\min},\\
\eta_0 & w_{\min}<W_r\le w_{\high},\\
\eta(w) & W_r>w_{\high}.
\end{cases}
\end{equation}

\paragraph{Error propagation.}
Define $\xi_R:=0$ and for $r=R-1,\dots,0$ set
\[
\xi_r\ :=\ \eta_r+(1+\eta_r)\xi_{r+1}.
\]
We claim that for all $r$,
\begin{equation}\label{eq:xi_claim_v2}
(1-\xi_r)W_r\ \le\ \widehat W_r\ \le\ (1+\xi_r)W_r.
\end{equation}
We prove this by backward induction on $r$.
The base case $r=R$ holds since $\widehat W_R=W_R$ and $\xi_R=0$.

Assume~\eqref{eq:xi_claim_v2} holds for $r+1$. Then
\[
\widehat W_r
=
|S\cap I_r|+2\widehat W_{r+1}
\le
|S\cap I_r|+2(1+\xi_{r+1})W_{r+1}.
\]
Using $2W_{r+1}=|R_r|+E_r$ where $|E_r|\le \eta_r W_r$ from~\eqref{eq:eta_r_def_v2}, and
$2W_{r+1}\le |R_r|+\eta_r W_r\le (1+\eta_r)W_r$, we get $2W_{r+1}\le (1+\eta_r)W_r$ and hence
$2\xi_{r+1}W_{r+1}\le (1+\eta_r)\xi_{r+1}W_r$.
Therefore
\[
\widehat W_r
\le
|S\cap I_r|+|R_r|+\eta_r W_r+(1+\eta_r)\xi_{r+1}W_r
=
W_r+\xi_r W_r
=
(1+\xi_r)W_r.
\]
A symmetric argument gives the lower bound:
\begin{align*}
\widehat W_r
&   =
|S\cap I_r|+2\widehat W_{r+1}
\ge
|S\cap I_r|+2(1-\xi_{r+1})W_{r+1}
=
|S\cap I_r|+|R_r|-E_r-2\xi_{r+1}W_{r+1}
\\
& \ge
W_r-\eta_r W_r-(1+\eta_r)\xi_{r+1}W_r
=
(1-\xi_r)W_r.
\end{align*}
    

This proves~\eqref{eq:xi_claim_v2} for all $r$, in particular
$\widehat W_0\in(1\pm \xi_0)W_0$.

Finally, since $1+\xi_r=(1+\eta_r)(1+\xi_{r+1})$, we have
\[
1+\xi_0=\prod_{r=0}^{R-1}(1+\eta_r)\ \le\ \exp\!\left(\sum_{r=0}^{R-1}\eta_r\right),
\]
hence $\xi_0\le \exp(\sum_r\eta_r)-1$.

\paragraph{Bounding $\sum_r\eta_r$.}
The number of rounds with $w_{\min}<W_r\le w_{\high}$ is $O(\log(2w_{\high}))$ since
$W_{r+1}\le (1+\eta_0)W_r/2\le (3/5)W_r$.
Thus $\sum_{w_{\min}<W_r\le w_{\high}}\eta_r\le O(\eta_0\log(2w_{\high}))\le \varepsilon/50$.

For rounds with $W_r>w_{\high}$, we have $\eta_r\le \eta(w_r)$ with $w_r\asymp W_r$ dyadic and
$W_{r+1}\le (1+\eta_r)W_r/2\le (5/8)W_r$. A geometric-series estimate gives
\[
\sum_{W_r>w_{\high}}\eta_r
\ \le\
O\!\left(\sqrt{\frac{k\log(e w_{\high}/k)+\log R}{w_{\high}}}\right)
\ \le\ \varepsilon/50
\]
by the definition of $w_{\high}$ (choosing $C$ large enough).
Therefore $\sum_r\eta_r\le \varepsilon/25$, and for $\varepsilon\le 1/4$ we get
$\xi_0\le e^{\varepsilon/25}-1\le \varepsilon$ (after increasing constants).
Hence $\widehat{\mathrm{wt}}_{T,\alpha}(x)=\widehat W_0\in (1\pm\varepsilon)W_0$ for all $x\in\calC$.}

\paragraph{Size bound.}
We have $|T|\le \sum_{r=0}^{R-1}|I_r|+|U_R|$ and $|U_R|\le 50$ on the good event.
Also $|I_{r,\mathrm{low}}|\le k w_{\low}$.

For each dyadic $w\in[w_{\low},w_{\high}]$, apply
\Cref{thm:universe_reduction} to the layer family
$\calF^{(r)}_{(w,2w]}$ with $\lambda = \theta_0 \eta_0^2$ to obtain
\[
|I_{r,w}|
\ \le\
\frac{k}{\eta_0^2}\cdot \poly\!\left(\log(kw/\eta_0)\right).
\]
Since $w\le w_{\high}$ throughout, we have $\log w\le \log w_{\high}$ and
\[
\log(k/\eta_0)=\log(k/\varepsilon)+O(\log\log w_{\high}),
\]
hence
\[
|I_{r,w}|
\ \le\
\frac{k}{\eta_0^2}\cdot \poly\!\left(\log(k/\varepsilon),\log\log n\right),
\]
using $\log w_{\high} = O(\log(k/\varepsilon)+\log\log n)$ by the definition of $w_{\high}$.

There are $O(\log(2w_{\high}))$ such dyadic values. Thus
\[
|I_r|
\ \le\
\frac{k}{\eta_0^2}\cdot \poly\!\left(\log(k/\varepsilon),\log\log n\right),
\]
since $\eta_0^{-2}=\Theta(\log^2(2w_{\high})/\varepsilon^2)$ and
$w_{\low}=\Theta(\eta_0^{-2}(\log(k/\varepsilon)+\log R))$.
Finally $R=\Theta(\log n)$ gives
\[
|T|
\ \le\
\frac{k\log n}{\varepsilon^2}\cdot \poly\!\left(\log(k/\varepsilon),\log\log n\right),
\]
as claimed.

This finishes the proof of the full sparsification theorem.
\end{proof}

\end{document}
